\AtBeginDocument{%
\begingroup\pagestyle{empty}\raggedleft\parindent0pt
\titulagemfront{}
\vfill
\clearpage

%% Créditos ------------------------------------------------------
\raggedright
\linha{copyright}{\copyrightlivro}
\linhalayout{edição brasileira©}{Circuito \ifdef{\ano}{\ano}{\the\year}}
\linha{tradução©}{\copyrighttraducao}
\linha{organização©}{\copyrightorganizacao}
\linha{introdução©}{\copyrightintroducao}
\linha{ilustração©}{\copyrightilustracao}\smallskip
\linha{título original}{\titulooriginal}
\linha{edição consultada}{\edicaoconsultada}
\linha{primeira edição}{\primeiraedicao}
\linha{agradecimentos}{\agradecimentos}
\linha{indicação}{\indicacao}
\linha{organização da coleção}{\organizador}
\linha{edição}{\edicao}
\linha{coedição}{\coedicao}
\linha{assistência editorial}{\assistencia}
\linha{preparação}{\preparacao}
\linha{revisão}{\revisao}
\linha{iconografia}{\iconografia}
\linha{capa}{\capa}
\linha{imagem da capa}{\imagemcapa}
\linha{imagens de capa e miolo}{\imagens}
\linha{ISBN}{\ISBN}
\linha{eISBN}{\eISBN}\smallskip
%\linhalayout{coleção}{\versal{Ataque Frontal}}
\linhalayout{conselho consultivo}{\tiny Amilcar Parker,\\
\qquad \qquad \qquad \qquad \quad \, Cecília Coimbra (\versal{TNM}/\versal{RJ} e \versal{UFF}),\\
\qquad \qquad \qquad \qquad \quad \, Edson Passetti (\versal{PUC}-\versal{SP}),\\
\qquad \qquad \qquad \qquad \quad \, Eduardo Sterzi (\versal{UNICAMP}),\\
\qquad \qquad \qquad \qquad \quad \, Heliana Conde (\versal{UERJ}),\\
\qquad \qquad \qquad \qquad \quad \, João da Mata (\versal{SOMA}),\\
\qquad \qquad \qquad \qquad \quad \, Jorge Sallum (Hedra),\\
\qquad \qquad \qquad \qquad \quad \, Margareth Rago (Unicamp),\\
\qquad \qquad \qquad \qquad \quad \, Priscila Vieira (\versal{UFPR}),\\
\qquad \qquad \qquad \qquad \quad \, Salvador Schavelzon (\versal{UNIFESP}),\\
\qquad \qquad \qquad \qquad \quad \, Thiago Rodrigues (\versal{UFF})}

\medskip
%\begingroup\tiny
%\ifdef{\conselho}{\conselho}{\relax}
%\par\endgroup\bigskip

\begingroup \tiny

\emph{Filosofia Black Bloc é um texto de ficção. Qualquer\\
alusão a fatos ou pessoas é meramente incidental.}\\

\textit{Grafia atualizada segundo o Acordo Ortográfico da Língua\\
Portuguesa de 1990, em vigor no Brasil desde 2009.}\\

%\begin{vplace}[1]
\begin{flushleft}
\hspace{10pt}Dados Internacionais de Catalogação na Publicação -- CIP
\_\_\_\_\_\_\_\_\_\_\_\_\_\_\_\_\_\_\_\_\_\_\_\_\_\_\_\_\_\_\_\_\_\_\_\_\_\_\_\_\_\_\_\_\_\_\_\_\_\_\_\_\_\_\_\_\_\_\_\_\_\_\_\_\_\_\_\_\_\_\_
\end{flushleft}
M977 \hspace{9pt}Correa, Murilo Duarte Costa\\
\hspace{20pt}\parbox{185pt}{Filosofia black bloc / Murilo Duarte Costa Correa. Apresentação de Bruno Cava. Prefácio de Peter Pál Pelbart. -- Rio de Janeiro: Circuito, 2018. (Coleção Ataque)}

\hspace{20pt}156 p.: Il.\\[6pt]

\hspace{20pt}ISBN 978-85-9582-XXX-X\\[6pt]

\hspace{20pt}\parbox{185pt}{1. Filosofia. 2. Filosofia política. 3. Movimento social. 4. Sistema político. 5. Controle social 6. Junho de 2013 I. Título.}

\begin{flushleft}
\hspace{20pt}CDU 329 \hspace{138pt}CDD 320
\_\_\_\_\_\_\_\_\_\_\_\_\_\_\_\_\_\_\_\_\_\_\_\_\_\_\_\_\_\_\_\_\_\_\_\_\_\_\_\_\_\_\_\_\_\_\_\_\_\_\_\_\_\_\_\_\_\_\_\_\_\_\_\_\_\_\_\_\_\_\_\\
\end{flushleft}
%\end{vplace}

\vfill\textit{Direitos reservados em l\'íngua\\ portuguesa somente para o Brasil}\\\medskip

%
\textsc{editora circuito ltda.}\\ 
Rua Visconde de Inhaúma, 134, grupo 1215 - Centro\\
20091-007, Rio de Janeiro-\textsc{rj}, Brasil\\
Telefone/Fax +55 21 2205 3236\\\smallskip
editoracircuito.com.br\\
\bigskip
Foi feito o depósito legal.\\\endgroup
\pagebreak\raggedleft
%% Front ---------------------------------------------------------
% Titulo
\titulagem

{\Large \autor \par\vspace{1.5ex}}
\ifdef{\apresentador}{{\small {\apresentador} (\textit{apresentação})} \par}{}
\ifdef{\prefacio}{{\small {\prefacio} (\textit{prefácio})} \par}{}
%\ifdef{\organizador}{{\small {\organizador} (\textit{organização})} \par}{}
\ifdef{\introdutor}{{\small {\introdutor} (\textit{introdução})} \par}{}
\ifdef{\tradutor}{{\small {\tradutor} (\textit{tradução})}\par}{}\vspace{1.5cm}

{{\footnotesize{} \ifdef{\numeroedicao}{\numeroedicao}{1}ª edição} \par}

%logos
\vfill
\ifdef{\logo}{\IfFileExists{\logo}{\hfill\includegraphics[width=3cm]{\logo}\hfill\logoum{}\\ São Paulo\_\the\year}}{\includegraphics[width=3cm]{CIRCUITO.jpg}\break{} Rio de Janeiro\_\the\year}
%\includegraphics[width=.4\textwidth,trim=0 0 25 0]{logo.jpg}\\\smallskip
\par\clearpage\endgroup
% Resumo -------------------------------------------------------
\begingroup \footnotesize \parindent0pt \parskip 5pt \thispagestyle{empty} \vspace*{.25\textheight}\mbox{} \vfill
\baselineskip=.92\baselineskip
\IfFileExists{PRETAS.tex}{\begin{itemize}
\item \textbf{Filosofia black bloc} \lipsum[1]
  
\item \textbf{Murilo Duarte Costa Correa} \lipsum[1]
\end{itemize}

}{% 
\ifdef{\resumo}{\resumo\par}{}
\ifdef{\sobreobra}{\sobreobra}{}
\ifdef{\sobreautor}{\mbox{}\vspace{4pt}\newline\sobreautor}{}
\ifdef{\sobretradutor}{\newline\sobretradutor}{\relax}
\ifdef{\sobreorganizador}{\vspace{4pt}\newline\sobreorganizador}{\relax}\par}
\thispagestyle{empty} \endgroup
\ifdefvoid{\sobreautor}{}{\pagebreak\ifodd\thepage\paginabranca\fi}
% Sumário -------------------------------------------------------

\sumario{}
\IfFileExists{INTRO.tex}{\include{INTRO}}

\IfFileExists{TEXTO.tex}{\mbox{}\include{TEXTO}}
%\part[{{\def\break{}\titulo}}]{\titulo}
} % fim do AtBeginDocument

% Finais -------------------------------------------------------
\AtEndDocument{%
  %\publicidade

\pagebreak\ifodd\thepage\paginabranca\fi

\ifdef{\imagemficha}{\IfFileExists{\imagemficha}{\includegraphics[width=.7\textwidth]{\imagemficha}\par}}{}

\mbox{}\vfill\small\thispagestyle{empty}
\begin{center}
\begin{minipage}{.8\textwidth}
\centering\tiny\noindent{}Adverte-se aos curiosos que se imprimiu este livro \ifdef{\grafica}{na gráfica \grafica}{em nossas oficinas}, 
em \today \ifdef{\papelmiolo}{em papel \papelmiolo}, em tipologia \tipopadrao{}, com diversos sofwares livres, 
entre eles, Lua\LaTeX, git \& ruby. \ifdef{\RevisionInfo{}}{\par(v.\,\RevisionInfo)}{}\par \begin{center}\normalsize\adforn{64}\end{center}
\end{minipage}
\end{center}
}
